% Môn: Vật lý
% Lớp: 11
% Chương: Chương I. Dao động
% Bài: L11C1 Bài 2. Mô tả dao động điều hòa · Xác định trạng thái của vật tại một thời điểm xác định
% ===== Câu 67 =====
% ID: L11C1B2-111-TLN
% Mức: TH
\dangbt{Xác định trạng thái của vật tại một thời điểm xác định}
\begin{ex}
% ID: L11C1B2-111-TLN
% Mức: TH
Một vật dao động theo phương trình x = 6cos(2πt) cm. Li độ của vật tại t = 1/ $3\,\text{s}$ theo đơn vị cm bằng bao nhiêu?
\shortans{-3}
\loigiai{
x(1/3) = 6cos(2π/3) = - $3\,\text{cm}$.
}
\end{ex}

% ===== Câu 69 =====
% ID: L11C1B2-112-DS
% Mức: NB
\begin{ex}
% ID: L11C1B2-112-DS
% Mức: NB
Một vật dao động điều hòa theo phương trình $x = 4\cos(2\pi t + \dfrac{\pi}{2})\ (\mathrm{cm})$ 
	Chọn Đúng hoặc Sai cho các phát biểu sau:
\choiceTF
{Tại $t = 0$, vật ở vị trí biên dương}
{\True Tại $t = 0$, vật qua vị trí cân bằng theo chiều âm}
{\True Vận tốc tại $t = 0$ bằng $-8\pi\ \mathrm{cm/s}$}
{\True Gia tốc tại $t = 0$ bằng $0\ \mathrm{cm/s^2}$}
\loigiai{
\begin{itemchoice}
\item (Sai) Tại $t = 0$, vật ở vị trí biên dương. (Vì): \omega \sin(\omega t + \varphi) = -4 \cdot 2\pi \cdot \sin(\dfrac{\pi}{2}) = -8\pi$, chuyển động theo chiều âm $\Rightarrow$ Đúng
\item (Đúng) Tại $t = 0$, vật qua vị trí cân bằng theo chiều âm. (Vì): $v(0) =
\item (Đúng) Vận tốc tại $t = 0$ bằng $-8\pi\ \mathrm{cm/s}$. (Vì): Vận tốc tại $t=0$ là $-8\pi\ \mathrm{cm/s}\Rightarrow$ Đúng
\item (Đúng) Gia tốc tại $t = 0$ bằng $0\ \mathrm{cm/s^2}$.
\end{itemchoice}
}
\end{ex}

% ===== Câu 70 =====
% ID: L11C1B2-113-TLN
% Mức: TH
\begin{ex}
% ID: L11C1B2-113-TLN
% Mức: TH
Một vật dao động theo phương trình x = 8cos(πt) cm. Li độ của vật tại t = $1\,\text{s}$ theo đơn vị cm bằng bao nhiêu?
\shortans{-8}
\loigiai{
x(1) = 8cosπ = - $8\,\text{cm}$.
}
\end{ex}

% ===== Câu 71 =====
% ID: L11C1B2-114-DS
% Mức: VD
\begin{ex}
% ID: L11C1B2-114-DS
% Mức: VD
Một chất điểm dao động điều hòa theo phương trình $x = 5\cos\left(10\pi t + \dfrac{\pi}{3}\right)$ (cm). Khảo sát trạng thái chuyển động của chất điểm:
\choiceTF
{\True Khi pha của dao động bằng $\pi$ rad thì li độ của vật bằng $-5$ cm.}
{Khi pha của dao động bằng $\dfrac{\pi}{2}$ rad thì li độ của vật bằng $2,5$ cm.}
{Tại thời điểm $t = 0,1$ s, li độ của chất điểm là $2,5$ cm.}
{\True Tại thời điểm ban đầu $t = 0$, li độ của chất điểm là $2,5$ cm.}
\loigiai{
\begin{itemchoice}
\item (Đúng) Khi pha của dao động bằng $\pi$ rad thì li độ của vật bằng $-5$ cm. (Vì): Mệnh đề thứ nhất đúng vì khi pha dao động $\Phi = \pi$ rad, ta thế vào phương trình thu được: $x = 5\cos(\pi) = -5$ cm
\item (Sai) Khi pha của dao động bằng $\dfrac{\pi}{2}$ rad thì li độ của vật bằng $2,5$ cm. (Vì): Mệnh đề thứ hai sai vì khi pha dao động $\Phi = \dfrac{\pi}{2}$ rad, ta có: $x = 5\cos\left(\dfrac{\pi}{2}\right) = 0$ cm
\item (Sai) Tại thời điểm $t = 0,1$ s, li độ của chất điểm là $2,5$ cm. (Vì): Mệnh đề thứ sáu sai vì thế $t = 0,1$ s vào phương trình ta được: $x = 5\cos\left(10\pi \cdot 0,1 + \dfrac{\pi}{3}\right) = 5\cos\left(\pi + \dfrac{\pi}{3}\right) = -2,5$ cm
\item (Đúng) Tại thời điểm ban đầu $t = 0$, li độ của chất điểm là $2,5$ cm. (Vì): Mệnh đề thứ tư đúng vì thế $t = 0$ s vào phương trình ta được: $x_0 = 5\cos\left(\dfrac{\pi}{3}\right) = 2,5$ cm
\end{itemchoice}
}
\end{ex}

% ===== Câu 72 =====
% ID: L11C1B2-115-DS
% Mức: VD
\begin{ex}
% ID: L11C1B2-115-DS
% Mức: VD
Một vật dao động điều hòa có đồ thị li độ – thời gian hình sin. Tại thời điểm $t = 0$, vật ở vị trí biên âm $x_0 = -4$ cm. Đường biểu diễn đi qua vị trí cân bằng lần thứ hai tại thời điểm mốc $t = 0,75$ s. Khảo sát các mệnh đề sau:
\choiceTF
{\True Biên độ dao động của vật bằng $4$ cm.}
{\True Kể từ thời điểm ban đầu, vật đi qua vị trí cân bằng lần thứ nhất tại thời điểm $t = 0,25$ s.}
{\True Chu kì dao động của vật là $1,0$ s.}
{Pha ban đầu của dao động bằng $0$ rad.}
\loigiai{
\begin{itemchoice}
\item (Đúng) Biên độ dao động của vật bằng $4$ cm. (Vì): Mệnh đề thứ nhất đúng vì vật xuất phát từ biên âm nên khoảng cách từ vị trí cân bằng tới biên là $A = 4$ cm
\item (Đúng) Kể từ thời điểm ban đầu, vật đi qua vị trí cân bằng lần thứ nhất tại thời điểm $t = 0,25$ s. (Vì): Mệnh đề thứ hai đúng vì thời gian vật đi từ biên âm đến vị trí cân bằng lần đầu tiên là một phần tư chu kì: $t_1 = \dfrac{T}{4}$
\item (Đúng) Chu kì dao động của vật là $1,0$ s. (Vì): Mệnh đề thứ ba đúng vì thời gian từ lúc xuất phát ở biên âm đến lúc qua vị trí cân bằng lần thứ hai là ba phần tư chu kì: $\dfrac{3T}{4} = 0,75 \text{ s} \Rightarrow T = 1,0$ s. Do đó $t_1 = \dfrac{1,0}{4} = 0,25$ s
\item (Sai) Pha ban đầu của dao động bằng $0$ rad. (Vì): Mệnh đề thứ tư sai vì ban đầu vật ở vị trí biên âm nên pha ban đầu phải bằng $\varphi = \pi$ rad
\end{itemchoice}
}
\end{ex}

% ===== Câu 73 =====
% ID: L11C1B2-116-TLN
% Mức: TH
\begin{ex}
% ID: L11C1B2-116-TLN
% Mức: TH
Một vật dao động theo phương trình x = 8cos(πt) cm. Li độ của vật tại t = $1\,\text{s}$ theo đơn vị cm bằng bao nhiêu?
\shortans{22.12}
\loigiai{
x(1) = 8cosπ = - $8\,\text{cm}$.
}
\end{ex}

% ===== Câu 74 =====
% ID: L11C1B2-117-DS
% Mức: VD
\begin{ex}
% ID: L11C1B2-117-DS
% Mức: VD
Cho dao động điều hòa có phương trình li độ $x = 2\cos\left(2\pi t - \dfrac{\pi}{4}\right)$ (cm). Xét trạng thái dao động của vật tại thời điểm $t = 0,5$ s:
\choiceTF
{\True Pha của dao động tại thời điểm đang xét là $\dfrac{3\pi}{4}$ rad.}
{\True Li độ của vật tại thời điểm đó có giá trị là $-\sqrt{2}$ cm.}
{Tại thời điểm đó, vật đang chuyển động theo chiều dương của trục tọa độ.}
{\True Gia tốc của vật tại thời điểm đang xét có giá trị dương.}
\loigiai{
\begin{itemchoice}
\item (Đúng) Pha của dao động tại thời điểm đang xét là $\dfrac{3\pi}{4}$ rad. (Vì): Mệnh đề thứ nhất đúng vì pha dao động tại $t = 0,5$ s là: $\Phi = 2\pi \cdot 0,5 - \dfrac{\pi}{4} = \dfrac{3\pi}{4}$ rad
\item (Đúng) Li độ của vật tại thời điểm đó có giá trị là $-\sqrt{2}$ cm. (Vì): Mệnh đề thứ hai đúng vì li độ của vật tại $t = 0,5$ s là: $x = 2\cos\left(\dfrac{3\pi}{4}\right) = -\sqrt{2}$ cm
\item (Sai) Tại thời điểm đó, vật đang chuyển động theo chiều dương của trục tọa độ. (Vì): Mệnh đề thứ ba sai vì pha dao động $\Phi = \dfrac{3\pi}{4}$ rad thuộc góc phần tư thứ hai, có $\sin\Phi \gt  0$ nên vận tốc mang giá trị âm ($v \lt  0$), nghĩa là vật chuyển động theo chiều âm
\item (Đúng) Gia tốc của vật tại thời điểm đang xét có giá trị dương. (Vì): Mệnh đề thứ tư đúng vì gia tốc $a = -\omega^2 x$. Do li độ $x = -\sqrt{2}$ cm mang giá trị âm nên gia tốc có giá trị dương
\end{itemchoice}
}
\end{ex}

% ===== Câu 75 =====
% ID: L11C1B2-118-TLN
% Mức: VD
\begin{ex}
% ID: L11C1B2-118-TLN
% Mức: VD
Một vật dao động điều hòa có phương trình $x=10\cos\left(2\pi t+\dfrac{\pi}{6}\right)$ cm. Tính li độ tại $t=0,5$ s.
\shortans{-8.66}
\loigiai{
x = 10 $\cos(2\pi\cdot$ 0,5 + $\dfrac{\pi}{6})$ = 10 $\cos(\pi$ + $\dfrac{\pi}{6})$. = 10 $\cdot\left(-\dfrac{\sqrt{3}}{2}\right)$ = -5 $\sqrt{3}\approx$ -8,66.
}
\end{ex}

% ===== Câu 80 =====
% ID: L11C1B2-119-DS
% Mức: VD
\begin{ex}
% ID: L11C1B2-119-DS
% Mức: VD
Một vật dao động điều hòa có phương trình li độ $x = 6\cos\left(4\pi t - \dfrac{\pi}{6}\right)$ (cm). Đánh giá tính đúng/sai của các nhận định dưới đây:
\choiceTF
{\True Tại thời điểm ban đầu $t = 0$, vật đang chuyển động theo chiều dương của trục tọa độ.}
{Li độ của vật ở thời điểm $t = 0,25$ s là $3\sqrt{3}$ cm.}
{\True Pha của dao động tại thời điểm $t = 1$ s là $\dfrac{23\pi}{6}$ rad.}
{\True Ở thời điểm $t = 0,5$ s, vật đi qua vị trí có tọa độ li độ $x = 3\sqrt{3}$ cm.}
\loigiai{
\begin{itemchoice}
\item (Đúng) Tại thời điểm ban đầu $t = 0$, vật đang chuyển động theo chiều dương của trục tọa độ. (Vì): Mệnh đề thứ nhất đúng vì pha ban đầu $\varphi = -\dfrac{\pi}{6} \lt  0$, dẫn đến vận tốc ban đầu của vật có giá trị dương ($v_0 \gt  0$), tức là vật chuyển động theo chiều dương
\item (Sai) Li độ của vật ở thời điểm $t = 0,25$ s là $3\sqrt{3}$ cm. (Vì): Mệnh đề thứ hai sai vì tại $t = 0,25$ s, li độ của vật là: $x = 6\cos\left(4\pi \cdot 0,25 - \dfrac{\pi}{6}\right) = 6\cos\left(\pi - \dfrac{\pi}{6}\right) = -3\sqrt{3}$ cm
\item (Đúng) Pha của dao động tại thời điểm $t = 1$ s là $\dfrac{23\pi}{6}$ rad. (Vì): Mệnh đề thứ ba đúng vì thế $t = 1$ s vào biểu thức pha dao động ta thu được: $\Phi = 4\pi \cdot 1 - \dfrac{\pi}{6} = \dfrac{23\pi}{6}$ rad
\item (Đúng) Ở thời điểm $t = 0,5$ s, vật đi qua vị trí có tọa độ li độ $x = 3\sqrt{3}$ cm. (Vì): Mệnh đề thứ tư đúng vì thế $t = 0,5$ s vào phương trình li độ ta thu được: $x = 6\cos\left(4\pi \cdot 0,5 - \dfrac{\pi}{6}\right) = 6\cos\left(2\pi - \dfrac{\pi}{6}\right) = 3\sqrt{3}$ cm
\end{itemchoice}
}
\end{ex}

% ===== Câu 85 =====
% ID: L11C1B2-120-DS
% Mức: VD
\begin{ex}
% ID: L11C1B2-120-DS
% Mức: VD
Một chất điểm dao động điều hòa theo phương trình $x = 10\cos\left(\pi t + \dfrac{\pi}{2}\right)$ (cm). Khảo sát chuyển động của chất điểm này:
\choiceTF
{\True Tại thời điểm ban đầu $t = 0$, chất điểm đi qua vị trí cân bằng theo chiều âm.}
{Tại thời điểm $t = 1$ s, li độ của chất điểm có giá trị là $10$ cm.}
{\True Quãng đường chất điểm đi được sau khi thực hiện xong một dao động toàn phần là $40$ cm.}
{Chất điểm mất đúng $0,5$ s để di chuyển từ vị trí cân bằng đến biên dương lần đầu tiên.}
\loigiai{
\begin{itemchoice}
\item (Đúng) Tại thời điểm ban đầu $t = 0$, chất điểm đi qua vị trí cân bằng theo chiều âm. (Vì): Mệnh đề thứ nhất đúng vì tại $t = 0$, li độ $x_0 = 10\cos\left(\dfrac{\pi}{2}\right) = 0$ và pha ban đầu $\varphi = \dfrac{\pi}{2} \gt  0$ nên vật chuyển động theo chiều âm qua vị trí cân bằng
\item (Sai) Tại thời điểm $t = 1$ s, li độ của chất điểm có giá trị là $10$ cm. (Vì): Mệnh đề thứ hai sai vì tại $t = 1$ s, li độ của chất điểm là: $x = 10\cos\left(\pi \cdot 1 + \dfrac{\pi}{2}\right) = 10\cos\left(\dfrac{3\pi}{2}\right) = 0$ cm
\item (Đúng) Quãng đường chất điểm đi được sau khi thực hiện xong một dao động toàn phần là $40$ cm. (Vì): Mệnh đề thứ ba đúng vì trong một dao động toàn phần (một chu kì), quãng đường vật đi được luôn bằng $4A = 4 \cdot 10 = 40$ cm
\item (Sai) Chất điểm mất đúng $0,5$ s để di chuyển từ vị trí cân bằng đến biên dương lần đầu tiên. (Vì): Mệnh đề thứ tư sai vì ban đầu vật ở vị trí cân bằng chuyển động theo chiều âm, vật phải đi đến biên âm (mất $0,5$ s), quay lại vị trí cân bằng (mất thêm $0,5$ s) rồi mới đến biên dương lần đầu tiên (tổng thời gian là $1,5$ s
\end{itemchoice}
}
\end{ex}

% ===== Câu 89 =====
% ID: L11C1B2-121-DS
% Mức: VDC
\begin{ex}
% ID: L11C1B2-121-DS
% Mức: VDC
Một vật dao động điều hòa được mô tả bằng phương trình li độ $x = 4\cos\left(2\pi t + \dfrac{\pi}{3}\right)$ (cm). Khảo sát các mốc thời gian chuyển động của vật:
\choiceTF
{\True Vật đạt li độ cực đại dương lần đầu tiên tại thời điểm $t = \dfrac{5}{6}$ s.}
{\True Vật đi qua vị trí cân bằng lần đầu tiên tại thời điểm $t = \dfrac{1}{12}$ s.}
{\True Khoảng thời gian ngắn nhất để vật di chuyển từ vị trí biên âm sang vị trí biên dương là $0,5$ s.}
{Tại thời điểm $t = 1$ s, vật đang ở vị trí cân bằng.}
\loigiai{
\begin{itemchoice}
\item (Đúng) Vật đạt li độ cực đại dương lần đầu tiên tại thời điểm $t = \dfrac{5}{6}$ s. (Vì): Mệnh đề thứ nhất đúng vì vật đạt biên dương khi pha dao động bằng $2k\pi$. Thời điểm đầu tiên $t \gt  0$ thỏa mãn: $2\pi t + \dfrac{\pi}{3} = 2\pi \Rightarrow t = \dfrac{5}{6}$ s
\item (Đúng) Vật đi qua vị trí cân bằng lần đầu tiên tại thời điểm $t = \dfrac{1}{12}$ s. (Vì): Mệnh đề thứ hai đúng vì vật qua vị trí cân bằng khi pha dao động bằng $\dfrac{\pi}{2} + k\pi$. Lần đầu tiên tương ứng với: $2\pi t + \dfrac{\pi}{3} = \dfrac{\pi}{2} \Rightarrow t = \dfrac{1}{12}$ s
\item (Đúng) Khoảng thời gian ngắn nhất để vật di chuyển từ vị trí biên âm sang vị trí biên dương là $0,5$ s. (Vì): Mệnh đề thứ ba đúng vì thời gian từ biên âm sang biên dương là nửa chu kì: $\Delta t = \dfrac{T}{2} = \dfrac{1}{2} = 0,5$ s
\item (Sai) Tại thời điểm $t = 1$ s, vật đang ở vị trí cân bằng. (Vì): o $T = \dfrac{2\pi}{2\pi} = 1$ s
\end{itemchoice}
}
\end{ex}

% ===== Câu 116 =====
% ID: L11C1B2-122-TN
% Mức: NB
\begin{ex}
% ID: L11C1B2-122-TN
% Mức: NB
Một chất điểm dao động dọc theo trục $Ox$. Phương trình dao động là $x = 2\cos(\pi t + \pi)$ (cm). Thời gian ngắn nhất vật đi từ lúc bắt đầu dao động đến lúc vật có li độ $x = \sqrt{3}$ là:
\choice
{$2,4\,\text{s}$}
{$1,2\,\text{s}$}
{\True $\dfrac{5}{6}$ s}
{$\dfrac{5}{12}$ s}
\loigiai{
Phương trình dao động của chất điểm là $x = 2\cos(\pi t + \pi)$ cm.
Ta có thể viết lại phương trình dưới dạng $x = -2\cos(\pi t)$ cm.
Tại thời điểm $t=0$, li độ của vật là $x(0) = -2\cos(0) = -2$ cm.
Vật bắt đầu dao động từ vị trí biên âm.
Ta cần tìm thời gian ngắn nhất để vật đi từ lúc bắt đầu dao động ($t=0$) đến lúc có li độ $x = \sqrt{3}$ cm.
Thay $x = \sqrt{3}$ vào phương trình dao động:
$\sqrt{3} = -2\cos(\pi t)\cos(\pi t) = -\dfrac{\sqrt{3}}{2}$
}
\end{ex}

% ===== Câu 117 =====
% ID: L11C1B2-123-TN
% Mức: NB
\begin{ex}
% ID: L11C1B2-123-TN
% Mức: NB
$\immini$ {Tại thời điểm $t = 0,5$ s, vận tốc của vật có độ lớn là:
\choice
{$0\ \mathrm{cm/s}$}
{\True $20\pi\ \mathrm{cm/s}$}
{$10\pi\ \mathrm{cm/s}$}
{$40\pi\ \mathrm{cm/s}$}
\loigiai{
Tại vị trí $t=0,5s$ thì đang ở Biên, do đó vận tốc bằng 0
}
\end{ex}

% ===== Câu 118 =====
% ID: L11C1B2-124-TN
% Mức: NB
\begin{ex}
% ID: L11C1B2-124-TN
% Mức: NB
Đồ thị biểu diễn li độ của một chất điểm dao động điều hòa như hình bên. 
 Tính quãng đường vật đi trong 3/4 chu kỳ.}{
 
	}
\choice
{$5\ \mathrm{cm}$}
{$10\ \mathrm{cm}$}
{\True $15\ \mathrm{cm}$}
{$20\ \mathrm{cm}$}
\loigiai{
$1$ chu kỳ vật đi $4A = 20$ cm 
 , $3/4$ chu kỳ: $S = 3A = 15$ cm.
}
\end{ex}

% ===== Câu 119 =====
% ID: L11C1B2-125-TN
% Mức: NB
\begin{ex}
% ID: L11C1B2-125-TN
% Mức: NB
Một chất điểm dao động điều hòa với phương trình li độ $x=2\cos\left(2\pi t+\dfrac{\pi}{2}\right)$ ($x$ tính bằng $\mathrm{~cm}, t$ tính bằng $s)$. Tại thời điểm $t=\dfrac{1}{4}\mathrm{~s}$ chất điểm có li độ bằng
\choice
{$2$ cm}
{$-\sqrt{3}$ cm}
{$\sqrt{3}$ cm}
{\True - $2$ cm}
\loigiai{
Thay $t=\dfrac{1}{4}$ s vào phương trình ta được $x=2\cos\left(2\pi\times\dfrac{1}{4}+\dfrac{\pi}{2}\right)=-2\mathrm{~cm}$
}
\end{ex}

% ===== Câu 120 =====
% ID: L11C1B2-126-TN
% Mức: NB
\begin{ex}
% ID: L11C1B2-126-TN
% Mức: NB
$\immini$ {Một vật dao động điều hòa theo phương trình $x = 6\cos\left(4\pi t\right)$ (cm). 
 Thời điểm đầu tiên vật đi qua vị trí $x = -3$ cm là:
\choice
{$t = \dfrac{1}{12}$ s}
{\True $t = \dfrac{1}{6}$ s}
{$t = \dfrac{1}{4}$ s}
{$t = \dfrac{1}{8}$ s}
\loigiai{
Ta có $x = 6\cos(4\pi t) = -3 \Rightarrow \cos(4\pi t) = -0,5$, $4\pi t = \dfrac{2\pi}{3} \Rightarrow t = \dfrac{1}{6}$ s (là nghiệm dương đầu tiên)
}
\end{ex}

% ===== Câu 121 =====
% ID: L11C1B2-127-TN
% Mức: NB
\begin{ex}
% ID: L11C1B2-127-TN
% Mức: NB
Chất điểm dao động điều hòa với phương trình $x=5\cos(4\pi$ t $+\frac{\pi}{3})$ (cm). Li độ của chất điểm khi pha dao động bằng $\frac{\pi}{3}$ là:
\choice
{– $2{,}5$ cm}
{$5$ cm}
{$0$ cm}
{\True $2{,}5$ cm}
\loigiai{
Khi pha dao động bằng $\frac{\pi}{3}$, li độ của chất điểm là $x=5\cos(\frac{\pi}{3})=2,5$ cm.
}
\end{ex}

% ===== Câu 122 =====
% ID: L11C1B2-128-TN
% Mức: NB
\begin{ex}
% ID: L11C1B2-128-TN
% Mức: NB
$\immini$ {Tại thời điểm $t = 1$ s, li độ của vật theo đồ thị là:
\choice
{\True $0\ \mathrm{cm}$}
{$20\ \mathrm{cm}$}
{$-20\ \mathrm{cm}$}
{$10\ \mathrm{cm}$}
\loigiai{
Tại $t = 1$ s, đồ thị cắt trục hoành , li độ $x = 0\ \mathrm{cm}$.
}
\end{ex}

% ===== Câu 123 =====
% ID: L11C1B2-129-TN
% Mức: NB
\begin{ex}
% ID: L11C1B2-129-TN
% Mức: NB
Một vật dao động điều hoà theo phương trình $x=4 \cos(7\pi t + \pi/6)$ cm. Vận tốc và gia tốc của vật ở thời điểm $t=2$ s là:
\choice
{14 $\pi$ cm/s và $-98\pi^2$ cm/s $^2$}
{\True $-14\pi$ cm/s và $-98\sqrt{3}\pi^2$ cm/s $^2$}
{$-14\pi\sqrt{3}$ cm/s và $98\pi^2$ cm/s $^2$}
{$14\,\text{cm}$ /s và $98\sqrt{3}\pi^2$ cm/s $^2$}
\loigiai{
$\omega = 7\pi$, $x = 4\cos(7\pi \cdot 2 + \pi/6) = 4\cos(14\pi + \pi/6) = 4\cos(\pi/6) = 2\sqrt{3}$.
$v = -\omega A \sin(\omega t + \varphi) = -28\pi \cdot \sin(\pi/6) = -14\pi$.
$a = -\omega^2 x = -(7\pi)^2 \cdot 2\sqrt{3} = -98\pi^2 \cdot \sqrt{3}$.
}
\end{ex}

% ===== Câu 124 =====
% ID: L11C1B2-130-TN
% Mức: NB
\begin{ex}
% ID: L11C1B2-130-TN
% Mức: NB
$\immini$ {Nếu tại $t = 0,5$ s vật có động năng cực đại, thì thế năng tại thời điểm này là:
\choice
{\True $0\ \mathrm{mJ}$}
{$100\ \mathrm{mJ}$}
{$50\ \mathrm{mJ}$}
{$20\ \mathrm{mJ}$}
\loigiai{
Tại $v = v_{\max}$ , $W_t = 0$ , thế năng bằng 0.
}
\end{ex}
